\documentclass{article}
\usepackage{xltxtra}
\usepackage{fontspec}
\usepackage{polyglossia}
\setmainlanguage{russian}
\setotherlanguage{english}
\setmainfont{Times New Roman}
\newfontfamily{\cyrillicfont}{Times New Roman}
\usepackage{amsmath}
\usepackage{physics}
\usepackage{tikz}
\usetikzlibrary{arrows.meta,calc}
\tikzset{
    >={Stealth[width'=2pt 0.1, length=5pt]}
}
\begin{document}
    \begin{tikzpicture}
        \draw (-1, 0) -- (1, 0);
        \foreach \i in {-1, -0.8, ..., 1 }
            \draw (\i, 0) -- (\i + 0.2, 0.2);
        \coordinate (O) at (-2, -4);
        \coordinate (mg) at ($(O) - (0, 2)$);
        \coordinate (q) at ($(O) -(1, 0)$);
        \draw (0,0) -- (O);
        \draw (O) circle (0.5);
        \draw[->] (O) -- node[right] {$\vec{mg}$} (mg);
        \draw[->] (O) -- node[above] {$\vec{F_{\text{к}}}$} (q);
        \draw[->] (O) -- node[right] {$\vec{F}$} ($(O) - (1, 2)$) coordinate (f);
        \draw[dashed] (mg) -- (f);
        \draw[dashed] (q) -- (f);
    \end{tikzpicture}
    
    Когда заряженное тело находиться в плоском конденсаторе на него действует сила Кулона:
    \[
        \vec{F_{\text{к}}} = q \cdot E
    \]
    где $q$ - заряд тела, а $E$ - напряженность электрического поля, которая равна 
    \[
        E = \frac{U}{d}
    \]
    где $U$ - напряжение на конденсаторе, $d$ - расстояние между пластинами конденсатора.

    Кроме этого на тело действует сила тяжести $\vec{mg}$. Равнодействующая этих сил 
    \[
        \vec{F} = \vec{F_{\text{к}}} + \vec{mg}
    \]
    Эта сила $\vec{F}$ растягивает нить на $\Delta x = 0.5\ \text{мм}$ Кроме того мы знаем коэффициент упругости нити $k = 100\ \frac{\text{Н}}{\text{м}}$

    Следовательно:
    \[
        F = \Delta x \cdot k
    \]
    Модуль силы Кулона будет равен:
    \[
        F_{\text{к}} = \sqrt{F^2 - (mg)^2}
    \]
    \[
        F_{\text{к}} = \sqrt{(\Delta x \cdot k)^2 - (mg)^2}
    \]
    С другой стороны сила Кулона будет равна:
    \[
        F_{\text{к}} = q \cdot \frac{U}{d}
    \]
    Отсюда:
    \[
        q \cdot \frac{U}{d} = \sqrt{(\Delta x \cdot k)^2 - (mg)^2}
    \]
    Выразим отсюда напряжение на конденсаторе:
    \[
        U = \frac{d \cdot \sqrt{(\Delta x \cdot k)^2 - (mg)^2}}{q}
    \]
    Подставим в формулу значения:
    \[
        U = \frac{5 \cdot 10^{-2}\ \text{м} \cdot \sqrt{(0.5 \cdot 10^{-3}\ \text{м} \cdot 100\ \frac{\text{Н}}{\text{м}})^2 - (3 \cdot 10^{-3}\ \text{кг} \cdot 9.8\ \frac{\text{м}}{\text{с}^2})^2}}{4 \cdot 10^{-7}\ \text{Кл}} \approx 5000 \ \text{В}
    \]
\end{document}